\section{Actual local branches of a general normal crossings carrier}
\label{sec:hironaka-general-nc-local-branches}

A globally self-intersecting hypersurface can have normal crossings
without admitting a global list of smooth members.  The local notion
below describes its carrier by actual coordinates.  Local smooth
branches and their complete divisor-ideal dictionary are then produced
from those coordinates.  No global labelling or descent of blowup
sequences is asserted.

\begin{definition}[A general normal crossings carrier]
\label{hir:general-normal-crossings-carrier}
\uses{hir:snc}
\lean{Hironaka.HironakaNormalCrossingsLocalBranches.IsNormalCrossingsCarrier}
\leanok
A closed set $H$ has normal crossings if, at every ambient point $x$,
there are centred holomorphic coordinates $\psi$ on a neighbourhood
and a set $S$ of distinct coordinate indices such that
\[
 y\in H\quad\Longleftrightarrow\quad
 \text{there exists }j\in S\text{ with }\psi_j(y)=0
\]
throughout their actual domain.  The set $S$ is allowed to be empty.
The definition imposes no global smoothness, connectedness, or
irreducibility on $H$.
\end{definition}

\begin{lemma}[The actual coordinate hyperplanes are SNC]
\label{hir:actual-coordinate-hyperplane-snc-list}
\uses{hir:snc,hir:snc-invariance}
\lean{Hironaka.HironakaNormalCrossingsLocalBranches.coordinateHyperplanes,Hironaka.HironakaNormalCrossingsLocalBranches.isSNCDivisor_coordinateHyperplanes}
\leanok
For a finite set $S$ of coordinate indices in $\mathbb C^n$, the list
of distinct actual hyperplanes $\{z_j=0\}$, $j\in S$, is SNC.  Their
traces in every open subset are SNC as well.
\end{lemma}
\begin{proof}
Each hyperplane is closed because a coordinate projection is continuous.
At any point $w$, translate the coordinates by $w$.  For a hyperplane
through $w$ its coordinate $w_j$ is zero, so the translated $j$th
coordinate vanishes exactly on that hyperplane.  The assigned singleton
coordinate sets are distinct because the indices in $S$ are distinct.
Translation is a genuine holomorphic equivalence with inverse translation
by $w$, and is centred at zero.  These are precisely the actual SNC
coordinate conditions.  Passing to an open subset preserves closedness
relative to that subset and restricts the same adapted coordinates.
The construction also applies to $n=0$, when $S$ and the list are empty.
\end{proof}

\begin{lemma}[The actual flattening map is a biholomorphism]
\label{hir:actual-flat-domain-biholomorphism}
\uses{hir:snc-invariance,hir:open-biholomorphism}
\lean{Hironaka.HironakaNormalCrossingsLocalBranches.flatDomainDiffeomorph}
\leanok
Let $\psi=\Phi\circ e$ be actual holomorphic flattening coordinates,
where $e$ is the preferred chart and $\Phi$ is a local holomorphic
equivalence whose source lies in the actual chart target.  The map
$\psi$ is a biholomorphism from its open actual domain onto
$\Phi$'s open target.
\end{lemma}
\begin{proof}
On the actual source define the map by $\Phi(e(y))$.  On the actual
target define its inverse by $e^{-1}(\Phi^{-1}(w))$.
The source containment ensures that this inverse is in the chart source
and in the flattening domain.  The two inverse equations are exactly
the actual chart and local holomorphic equivalence inverse equations.
Both maps are holomorphic on their specified open domains, being
compositions of holomorphic chart maps and holomorphic equivalence
maps.  Complex holomorphic maps there are smooth of every order.
Thus these actual mutually inverse maps give the stated biholomorphism.
\end{proof}

\begin{theorem}[Actual smooth coordinate branches]
\label{hir:general-nc-actual-local-snc-branches}
\uses{hir:general-normal-crossings-carrier,hir:actual-coordinate-hyperplane-snc-list,hir:actual-flat-domain-biholomorphism}
\lean{Hironaka.HironakaNormalCrossingsLocalBranches.exists_actual_snc_branches}
\leanok
For every point $x$ of a manifold with a normal crossings carrier $H$,
there is an actual open neighbourhood $U$ and a finite SNC list $B$
of closed smooth hypersurfaces in $U$, such that
\[
 H\cap U=\bigcup_{C\in B}C.
\]
The list is the inverse image of distinct coordinate hyperplanes under
the actual coordinate biholomorphism.  In particular it is produced
from the geometric coordinate condition rather than supplied as input.
\end{theorem}
\begin{proof}
Choose the actual coordinates and $S$ from the definition at $x$,
and let $U$ be their actual flattening domain.  It is open and contains
$x$.  For each $j\in S$ take the actual set
$C_j=\{y\in U:\psi_j(y)=0\}$.
Under the preceding actual biholomorphism it is the inverse image of
the corresponding coordinate hyperplane in the open target.
The coordinate-hyperplane lemma and biholomorphic invariance produce
the SNC condition on the whole of $U$.  Each branch is relatively
closed and smooth.  The defining coordinate equation gives the stated
carrier equality at every point of $U$.  Empty $S$ gives the empty
branch list and an empty carrier on $U$.
\end{proof}

\begin{lemma}[The actual original divisor and exponents on an open subset]
\label{hir:actual-divisor-open-germ-exponent-dictionary}
\uses{hir:divisor-exponent,hir:open-restriction-ideal,hir:open-restriction-marked,hir:open-biholomorphism}
\lean{Hironaka.HironakaNormalCrossingsLocalBranches.restrictDivisor,Hironaka.HironakaNormalCrossingsLocalBranches.ideal_restrictDivisor,Hironaka.HironakaNormalCrossingsLocalBranches.support_restrictDivisor,Hironaka.HironakaNormalCrossingsLocalBranches.mem_restrictDivisor_iff,Hironaka.HironakaNormalCrossingsLocalBranches.divExp_restrictDivisor_eq}
\leanok
For an actual effective divisor $D$ and an open subset $U$, its actual
restriction $D_U$ has ideal $I_{D_U}=\iota^*I_D$ and support
$\iota^{-1}\operatorname{Supp}D$, where $\iota:U\to M$ is inclusion.
For every point $y\in U$ and every germ represented by a function $g$
on $U$, membership in $I_{D_U,y}$ is precisely membership of the
extension of $g$ in $I_{D,y}$.  For every carrier $C\subset U$ and
$y\in C$, its genuine divisorial exponent satisfies
\[
 e_C(I_{D_U},y)=e_{\iota(C)}(I_D,\iota(y)).
\]
\end{lemma}
\begin{proof}
An open inclusion is holomorphic and has bijective differential at
every point.  Pullback therefore preserves actual invertibility,
giving the genuine effective divisor $D_U$.  The coherent pullback
zero-locus identity gives the asserted actual support.

For complete stalk membership, a function on $U$ may be extended by
zero outside $U$: near an original point of $U$ this is exactly its
holomorphic germ.  Finite holomorphic linear combinations of pulled-back
original ideal germs extend to the corresponding original combinations,
and every original germ restricts back.  Coherence identifies these
spans with the complete original ideal stalk, proving the equivalence
for every $g$.

Actual marked cosupports restrict under open inclusion, for every mark.
The inclusion identifies neighbourhood filters at $y$ with their
restriction from $M$.  Since the original image of $C$ and $C$ have
exactly matching membership at points of $U$, divisibility near $y$
along $C$ is equivalent to divisibility near $\iota(y)$ along
$\iota(C)$ at every mark.  Both actual invertible divisor ideals have
nonzero germs in each stalk, so their genuine finite exponents are
determined by these identical divisibility thresholds.  This proves
the claimed original exponent equality.  The image carrier need not
be globally closed in $M$; the argument uses its actual germ at a point
of the open subset.
\end{proof}

\begin{theorem}[The full local divisor dictionary for general normal crossings]
\label{hir:general-nc-original-divisor-local-product-dictionary}
\uses{hir:general-nc-actual-local-snc-branches,hir:actual-divisor-open-germ-exponent-dictionary,hir:locally-finite-original-component-product}
\lean{Hironaka.HironakaNormalCrossingsLocalBranches.exists_original_local_component_dictionary}
\leanok
Suppose the actual effective divisor $D$ is supported on a general
normal crossings carrier $H$.  At every original point $x$, produce
the open subset $U$ and actual SNC branch list $B$ above.  The genuine
component family $\mathcal C(B)$ is locally finite, and there is one
function $a$ on this whole family such that
\[
 a(C)=e_{\iota(C)}(I_D,\iota(y))\qquad(y\in C).
\]
Every point of $U$ has an open neighbourhood $V$ in $U$ and a finite
list of distinct actual carriers $A=((C_i,a(C_i)))$ with positive
coefficients, satisfying at every $y\in V$ the complete original
stalk dictionary
\[
 g\in(I_{D_U})_y
 \quad\Longleftrightarrow\quad
 g\in\left(\prod_i I_{C_i}^{a(C_i)}\right)_y
 \qquad(g:U\to\mathbb C).
\]
The left-hand side also equals membership of the extension of $g$
in the actual original ideal of $D$.  No compactness or global smooth
boundary list is assumed.
\end{theorem}
\begin{proof}
Produce the actual SNC list $B$ from the original coordinate condition,
and construct the actual restricted divisor $D_U$.
Its support is contained in $\iota^{-1}H=\bigcup B$ by original support
containment and the actual carrier equality.  Apply the genuine locally
finite SNC divisor-product theorem to this actual divisor and actual
branch list.  It produces $a$, local finiteness, and the positive finite
lists with the complete stalk equivalence.  The preceding exponent
transport identifies every produced coefficient with the exponent of
the original divisor on the image of the genuine carrier, at every
original carrier point.  The complete open-stalk equivalence supplies
the final original ideal interpretation.  All choices are produced
from the original geometric normal crossings condition and the actual
divisor; no factor or coefficient dictionary is an incoming premise.
\end{proof}

These are genuine local branches and actual divisor-ideal identities.
For a global normal crossings boundary whose branches have monodromy
or self-intersections, compatibility and descent of canonical blowup
sequences still require a separate proof.  The local dictionary does
not assert that global resolution theorem or close that remaining G9
obligation.
